\documentclass[10pt]{article}

\usepackage{graphicx}
\usepackage{comment}
\usepackage{wrapfig}
\usepackage{amsmath}
\usepackage{amssymb}
\usepackage{subfigure}
\usepackage{caption}
\usepackage{subcaption}
\usepackage{siunitx}
\usepackage{enumitem}
\usepackage{xcolor}
\usepackage{url}
\usepackage{fontspec}
\usepackage[backend=bibtex,bibstyle=mla]{biblatex}
\usepackage{filecontents}
\addbibresource{references.bib}
\setmainfont{Latin Modern Roman}
%\bibliographystyle{amsplain}
%\bibliography{references}


\title{Studio dello stato attivo del blazar Mrk421 con ASTRI Mini-Array​}

\author{Giacomo Spaccini}
\date{17 Dicembre 2026}

\begin{document}


\maketitle

\section{Introduzione}

\newpage

\section{Very High Energy AstroPhysics}

L'astrofisica contemporanea si divide in molte branche, tutte connesse dallo studio di tutti quei messaggeri cosmici che consentono di osservare gli eventi e processi che caratterizzano l'Universo.
\\
Tra i vari approcci esistenti allo studio del cosmo, quello legato al lavoro mostrato in questa tesi è l'Astrofisica delle Altissime Energie (in inglese \textit{Very High Energy AstroPhysics} o, in breve, VHEAP), la quale si occupa prevalentemente dello studio dei fotoni gamma che appartengono, per l'appunto, alla fascia delle altissime energie, ovvero da circa $ 30 \hspace{3 pt} \unit{GeV}$ fino al centinaio di $ \unit{TeV} $ \cite{GRA2018}.
\\

\subsection{Introduzione all'astrofisica delle altissime energie}

I fotoni sono da sempre stati i messaggeri privilegiati nello studio dell'Universo, sin dagli albori dello studio del cielo attraverso la luce visibile sino ai giorni d'oggi, dove lo spettro studiato varia in energia dai $\unit{keV}$ dello studio della radioastronomia e radioastrofisica fino ai $\unit{PeV}$ tipici dell'astrofisica delle energie estreme.
\\
Uno spettro che varia per oltre 12 ordini di grandezza, ed ogni banda spettrale corrisponde a certi tipi di informazioni estraibili sull'Universo e differenti sistemi di rivelazione.

Il primo ambito è quello delle onde radio, con lunghezze d'onda che variano tra i $ 10^{-3} \hspace{3 pt} \unit{m} $ ed i $ 10^2 \hspace{3 pt} \unit{m} $, che consentono di studiare sia sorgenti interne al sistema solare, come Giove o lo stesso Sole, che sorgenti più lontane, come ad esempio quasar, pulsar o nuclei galattici attivi (in inglese \textit{Active Galactic Nuclei}, in breve AGN) fino a consentire di rilevare la radiazione cosmica di fondo, ovvero la radiazione residua attribuita al Big Bang.
\\
Si passa poi all'infrarosso, la banda spettrale tipica delle emissioni di corpi con temperature che arrivano fino alle migliaia di kelvin e la cui radiazione tende a superare ostacoli come le nubi di polvere interstellare e consente di studiare sorgenti che sarebbero oscurate in altre bande.
\\
Diminuendo ancora la lunghezza d'onda ed arrivando tra gli $800 \hspace{3 pt} \unit{nm}$ ed i $300 \hspace{3 pt} \unit{nm}$ circa si giunge allo spettro visibile, la banda spettrale più a lungo studiata e che è dominata dall'emissione di corpi con temperature più alte, che arrivano alle decine di migliaia di kelvin, come ad esempio le stelle.

Come per le onde radio, anche la radiazione visibile è tranquillamente studiabile da terra, dato che l'atmosfera è abbondantemente trasparente a queste frequenze.
\\
Aumentando ulteriormente l'energia della radiazione osservata si incontrano le emissioni che ricadono nell'ultravioletto e, soprattutto, nei raggi X, radiazioni emesse prevalentemente da fenomeni molto intensi e corpi dalle temperature che superano i milioni di kelvin.

Nonostante la difficoltà nell'osservare questi fotoni per via dell'opacità dell'atmosfera, il che spesso richiede di avere rilevatori in orbita oltre la ionosfera o persino sonde spaziali, lo studio dei raggi X di origine extraterrestre ha consentito di analizzare processi non percepibili nell'ottico e popolando il cielo di tantissime sorgenti altrimenti invisibili.
\\
Infine si raggiunge lo spettro dei fotoni gamma, i fotoni più energetici
\\
\begin{comment}
L'interesse scientifico per questa categoria di fotoni è dovuto ad una serie di ragioni.
\\
In primo luogo, anche se l'atmosfera terrestre non è trasparente per nessuna radiazione gamma, questi fotoni sono sufficientemente energetici da poter essere studiati da terra, anche se indirettamente attraverso cascate atmosferiche o flash Cherenkov.
\\
Inoltre, lo studio dei fotoni gamma consente di osservare particolari ed estremi fenomeni astrofisici, come jet ultra-relativistici che sfuggono da buchi neri supermassivi, resti di supernova o \textit{gamma-ray bursts} provenienti da fonti come le blazar.
\\
Infine, molti degli eventi che producono raggi gamma generano anche raggi cosmici \cite{GRA2018}, particelle cariche dall'ampio spettro sia di flusso che di energia.
\\
Questi eventi danno luogo ad una emissione detta "non termica" (\textit{non-thermal emission}), ovvero l'emissione di energia da parte di un corpo che non si trova in equilibrio termodinamico e che conseguentemente non segue lo spettro di emissione di corpo nero, e l'origine di questi fotoni spesso si trova nella conversione di energia gravitazionale in radiazione elettromagnetica \cite{VHEGA}.
\end{comment}
\\

\newpage

\section{Telescopi Cherenkov ed ASTRI}

\newpage

\section{Panoramica di Markarian421}

\newpage

\section{Analisi}

\newpage

\section{Conclusioni}

\newpage

\section{Bibliografia}
\printbibliography

\end{document}
